Several systems combine embedded sensing with closed-loop control of soft or continuum arms. Hofer et al.\ placed one camera inside each of the three fabric bellows of a pneumatic spherical arm, observed a painted interior pattern under internal illumination, and regulated the two joint angles in closed loop with about $1^\circ$ accuracy at 30~Hz; the actuation redundancy of the arm was not exploited \cite{hoferVisionBasedSensingApproach2021}. Thuruthel et al.\ learned closed-loop inverse kinematics with end-effector feedback for a redundant tendon-driven continuum manipulator, with and without external loads \cite{thuruthelLearningClosedLoop2017}. Han et al.\ tracked an internal anchor marker with a single onboard camera to infer the full-body deformation and contact force of soft modules \cite{hanAnchoringMorphologicalRepresentations2025}, and Yoo et al.\ estimated the 3D shape of a pneumatic actuator from an internal camera with a model trained in simulation \cite{yooZeroShotSimtoRealTransfer2023}; neither closed a control loop on the estimate. Lu et al.\ servoed the 3D shape of continuum and soft robots with adaptive online learning, using fiber Bragg grating sensing \cite{luAdaptiveOnlineLearning2024}. The present work differs in two respects: it distributes one internal camera per section along a multi-section tendon-driven arm, without internal markers, patterns, or dedicated illumination, and it closes the loop on the camera-estimated pose with controllers that resolve the actuation redundancy explicitly, either in the null space of a model-based Jacobian or through a secondary objective in the training of a learned inverse kinematics. To the best of our knowledge, this combination has not been reported.
